\chapter{Лабораторная работа №1. Huawei VRP и основы конфигурирования}
\label{ch:chapter 1}

\section{Цели}

Лабораторная работа помогает получить практические навыки по изучению
следующих тем:

\begin{itemize}

\item Режимы командной строки, способы входа и выхода из режимов командной строки

\item Стандартные команды

\item Использование интерактивной справки командной строки

\item Отмена действия команды

\item Использование сочетаний клавиш в командной строке

\end{itemize}

\section{Топология сети}

Ниже представлена топология сети для данной работы.

\includegraphics[width=0.5\textwidth]{lb1_1.png} 
\captionof{figure}{Топология сети для лабораторной работы №1}

\section{План работы}

\begin{itemize}

\item 1. Выполнение базовых настроек, включая настройку имени устройства и IP-адреса
        интерфейса маршрутизатора.

\item 2. Сохранение конфигурации.

\item 3. Перезагрузка устройства.

\end{itemize}

\section{Процедура конфигурирования}

\textbf{Шаг 1.} Войдем в CLI маршрутизатора через консольный порт.


\includegraphics[width=0.5\textwidth]{lb1_2.png} 

\textbf{Шаг 2.} Выведем основную информацию об устройстве:

\includegraphics[width=0.5\textwidth]{lb1_3.png} 

\textbf{Шаг 3.} Настроим основные параметры устройства.

Изменим имя маршрутизатора на Datacom-Router:

\includegraphics[width=0.5\textwidth]{lb1_4.png} 

Войдем в режим интерфейса и настроим IP-адрес устройства:

\includegraphics[width=0.5\textwidth]{lb1_5.png} 

Теперь отменим настройку IP-адреса на интерфейсе GigabitEthernet 0/0/1 с помощью команды \textbf{undo}, т.к. IP-адрес предназначается для GigabitEthernet 0/0/2:

\includegraphics[width=0.5\textwidth]{lb1_6.png} 

Выведем текущую конфигурацию устройства с помощью команды \textbf{display current-configuration}

\includegraphics[width=0.5\textwidth]{lb1_7.png} 

\textbf{Шаг 4.} Сохраним текущую конфигурацию устройства:

Для этого сначала вернемся в режим пользователя командой \textbf{return}

\includegraphics[width=0.5\textwidth]{lb1_8.png} 

Её отличие от \textbf{quit} в том, что она позволяет сразу вернуться в режим пользователя, а не на шаг назад.

Сохраним конфигурацию при помощи команды \textbf{save}

\includegraphics[width=0.5\textwidth]{lb1_9.png} 

Как можно заметить, конфигурация сохраняется в файл, что удобно, если мы захотим использовать её для другого устройства, к примеру.

Теперь сравним текущую конфигурацию с конфигурацией в файле загрузки при помощи команды \textbf{compare}:

\includegraphics[width=0.5\textwidth]{lb1_10.png} 

\textbf{Шаг 5.} Выполнение операций в файловой системе.

Сначала выведем список всех файлов в текущем каталоге с помощью команды \textbf{dir}, как можно заметить, в маршрутизаторе Huawei можно работать с файлами, как в Linux, 
что является довольно удобным подходом

\includegraphics[width=0.5\textwidth]{lb1_10.png} 

Сохраним текущую конфигурцию в файл \textbf{test.cfg}:

\includegraphics[width=0.5\textwidth]{lb1_12.png}

Теперь проверим, создался ли он:

\includegraphics[width=0.5\textwidth]{lb1_13.png} 

Как можно заметить, он создался.

Зададим этот файл в качестве файла конфигурации. Для этого воспользуемся командой \textbf{startup saved-configuration}:

\includegraphics[width=0.5\textwidth]{lb1_14.png} 

Выведем информацию о файле конфигурации загрузки, с учетом его изменения:

\includegraphics[width=0.5\textwidth]{lb1_15.png} 

Удалим файл конфигурации:

\includegraphics[width=0.5\textwidth]{lb1_16.png} 

\textbf{Шаг 6.} Перезагрузим устройство.

Сделать это можно с помощью команды \textbf{reboot}

\includegraphics[width=0.5\textwidth]{lb1_17.png} 

\section{Вопросы}

На шаге 5 команда reset saved-configuration выполняется для удаления
конфигурации. Почему после перезапуска устройства конфигурация
сохраняется? - Дело в том, что мы удаляем файл test.cfg, но как видно на скриншоте на 5м шаге, это файл конфигурации при следующем запуске, а значит ничего и не поменялось.


